\SourceLeaf{289}%
\begin{center}
\begin{tabular}{lr}
Chasseral-Suchet&\(+0''04536\)\\
Chasseral-Berra&\(-0{,}00966\)\\
Suchet-Chasseral&\(+0{,}06221\)\\
Suchet-Berra&\(+0{,}01014\)\\
Berra-Suchet&\(-0{,}04717\)\\
Berra-Chasseral&\(-0{,}06039\)
\end{tabular}
\end{center}
also auch hier ohne Einfluss auf die Rechnung, die in dem angeführten Werke auf Zehntel der Secunde geführt ist.
\Article{15}
Die in den Artt. 12 und 13 behandelte Aufgabe ist zwar durch die gegebenen Vorschriften mit einer für die Anwendung überflüssig ausreichenden Genauigkeit aufgelöset; indessen ist es doch der Mühe werth, und zur gleichmässigen Vollendung einer in der Folge mitzutheilenden Untersuchung sogar nothwendig, für einen speciellen Fall die Genauigkeit noch um eine Ordnung weiter zu treiben: dieser specielle Fall steht unter der Bedingung, dass die Linie \(N\) in einem zwischen \(F\) und \(G\) liegenden Punkte \(H\) den Normalparallelkreis treffe. Es ist in diesem Falle vortheilhafter, den Anfangspunkt der \(x\), nicht wie oben in \(F\), sondern in \(H\) zu setzen, wodurch bewirkt wird, dass bei der Entwicklung von \(l\) und \(m\) in nach Potenzen von \(x\) fortschreitende Reihen in der erstern das erste und zweite Glied, in der andern das zweite und dritte ausfallen, oder dass sie folgende Form haben:
\[
\begin{aligned}
l&=\lambda xx+\lambda'x^3+\text{u.\,s.\,w.}\\
m&=1+\mu x^3+\mu'x^4+\text{u.\,s.\,w.}
\end{aligned}
\]
\indent Für unsern Zweck wird von den Coëfficienten in diesen Reihen nur der eine \(\lambda\) erforderlich sein, wofür sich aus der im 9. Art. für \(\log m\) gegebenen Formel verbunden mit den Entwicklungen des 13 Art. leicht folgender Ausdruck ableiten lässt:
\[
\lambda=-\frac{2ee\cos P\sin P\sin\chi\cos\chi^2}{\cos\varphi\cos\theta}
\]
in welcher \(e,P,\varphi,\theta\) ihre oben erklärten Bedeutungen behalten, und für \(\chi\) das in dem Punkte \(H\) Statt findende Azimuth des Bogens \(N\) zu setzen ist.

Werden obige Reihen bei der Integration der Gleichungen
